\section{结论}
\label{sec:conclusion}

本章把四问的答案汇总到一处。摘要说明本文做了什么，表~\ref{tab:overview} 给出四问各自交付的那一个方案，第~\ref{sec:evaluation} 章评价方法本身的好坏与可推广性；本章回答的是另一个问题——这四问的结论分别是什么、能拿去做什么、在什么条件下成立。本章只汇总第~\ref{sec:validation} 章复核过的结论，不新增计算，也不新增数字。

\subsection{四问的主要结论}

\textbf{问题一给出的是运力边界与最省架次的组批方案。}在允许同一服务区混用机型的可行域上，\mtQoneTrips\ 架次可覆盖全部 80 个货箱，总能耗 \mtQoneEnergy\,kWh、累计作业时间 \mtQoneWorkTimeH\,h；相较限定「一个服务区只用一个机型」时的 $63.180$\,kWh，混用机型在架次数不变的前提下节能 $6.24\%$。该结论另由三条互不相同的精确路径逐行比对，结果一致（表~\ref{tab:q1ilp}）。与之配套的是一条必须一并读的边界：返航安全余量超过 $0.3513$ 后本方案整体失效，而率先失去可达性的\textbf{不是}载重最大的 C 型、而是航程最短的 A 型——$\rho>0.2992$ 时 A 型空载亦无法往返最不利的 S008（5.3 节）。

\textbf{问题二给出的是资源受限下按时交付的 24 架次调度方案。}8 架运输无人机与 14 组共享电池全部投入周转，\mtQtwoRecTrips\ 个架次在 \mtQtwoMakespan\,s 内完成全部 80 个货箱的交付，加权时延与硬时限违约均为 $0$，总能耗 \mtQtwoEnergy\,kWh（表~\ref{tab:hardcheck}）。调度结构本身是效率的一等来源：在组批不变的条件下仅调整派工顺序，完成时间即下降 $31.9\%$。该方案附带一个不应被略过的条件——两层小规模精确验证中 \mtQtwoExactProved\ 例取得了最优性证书、\mtQtwoExactTimeout\ 例超时后如实以上界报告，而这些算例只含 3 个服务区、机队资源竞争几乎不存在，故它们的作用是对齐组批、能耗、路径与时序的逻辑，\textbf{不能}外推为本方案在大规模机队竞争下也已取得最优性证书（6.4 节）。

\textbf{问题三的结论是一条权衡：连续通信的硬约束会把问题二的高效方案判为不可执行，而消除中断要付出完工时刻的代价。}沿问题二的 \mtQtwoRecTrips\ 条三维轨迹逐时刻判定，只有 \mtQthreeDirect\ 的时间可直连、\mtQthreeRelayFrac\ 处于通信盲区（表~\ref{tab:commcheck}）。把运输开始时刻与错峰、悬停站址与离地高度、中继架次调度放进同一优化问题后，\mtQthreeStations\ 个悬停站、\mtQthreeRelays\ 次中继出动（共 \mtQthreeRelayVisits\ 次悬停站服务）以 \mtQthreeRelayEnergy\,kWh 把中断时间占比降至 \mtQthreeGap。代价是联合完工时刻由 \mtQtwoMakespan\,s 推迟到 \mtQthreeJointDone\,s（$+2.1\%$），而迟到箱数与加权迟到仍为 $0$、最小硬时限余量 \mtQthreeMinMarginFine\,s（表~\ref{tab:q3chain}）。还有一条只在本方案成立的观察：六站离地高度落在 $65\sim185$\,m、没有任何一站触及题给的 $300$\,m 上限，中继能量占用比最高仅 $0.740$，说明限制保障能力的是时序结构与站址几何而不是能量——换站址、加大服务窗口或放宽并发都会改变这一占用比，该观察不外推。

\textbf{问题四给出的不是一张时刻表，而是一条关于约束的发现。}按「同一架次访问的服务区不可拆分」把 15 个服务区绑定为 \mtQfourUnits\ 个原子任务单元后，$K=2$ 与 $K=3$ 分别有 \mtQfourPartKTwo\ 与 \mtQfourPartKThree\ 种分区，逐行枚举即得全局最优；按题面「资源不得跨组调配」复核，两类的库存可行分区数均为 \mtQfourFeasible（表~\ref{tab:q4check}），即「资源不得跨组调配」与题给库存在本算例下不相容。枚举进一步厘清了成因：八类资源各自的最小缺口在两 $K$ 下全部为 $0$，单独看任何一类都存在一个不超库存的分区，故不可行来自八类资源无法在同一分区内同时取到各自最小值的\textbf{联合约束}，而不是任何单一类别的短缺（表~\ref{tab:q4-min}）。其中中继无人机离临界最近：每个含失效区间的组都须自备通信保障，故需要中继的组数下界为 $K-1$，$K=3$ 时恰为 $2$、等于库存，余量为 $0$。推荐分区上的缺口是「先省资源、再求均衡」这一目标的代价：$K=2$ 缺 \mtQfourGapSum\ 件、$K=3$ 缺 \mtQfourMinPatchKThree\ 件。

\subsection{可信度与适用边界}

上述结论的可信度来自四条互不相同的复核路径（表~\ref{tab:overview}）：复核脚本\textbf{只读}四问的结果表重算硬约束，不复用求解器的任何中间变量，故它能排除程序内部状态与结果表之间的串号、漏检与自我印证。但该脚本与求解器仍共享公共物理模块，因此它\textbf{不能}排除公共物理模型本身的口径错误——后者只能靠 DEM 口径对照与原始数据判读（第~\ref{sec:dem-caliber} 节）。这条边界不在结论层放宽。

以下五条界定本文结论的适用范围，逐条见第 10.2 节。（1）问题二的推荐方案只在小规模算例上取得了最优性证书，其全局最优性在大规模机队竞争的情形下未被覆盖。（2）放开物理指派后得到的可行上界违反 9 项硬时限，只能作搜索方向上的佐证，其最优性间隙如实报出。（3）遮挡判定是二值的几何视线判据，未含第一菲涅尔区净空、绕射与多径，故直连时间占比应读作偏乐观的上界。（4）全文未建模运输无人机之间的最小安全间隔与高度层分配，冲突只在「同一资源不被两个任务同时占用」的意义上被排除。（5）问题四的否定结论只在本文 \mtQfourUnits\ 个原子单元的划分方式与题给库存下成立，换一种划分须重新核算。

\subsection{面向救援调度的建议}

把四问的结论落到调度决策上，本文给出四条建议，按可落地程度排序。

\textbf{第一，首选不分区、按单组执行。}这是全篇唯一无需任何补配即可交付四问的方案：以问题三的运输—中继联合调度整体执行，其资源需求不超过现有库存（8 架运输无人机与 14 组共享电池全部投入周转，见 6.3 节）。分区虽然便于并行指挥，却在题面「资源不得跨组调配」的规则下不可行。

\textbf{第二，若必须分区，则先补配再分区。}按推荐分区计，$K=2$ 缺 \mtQfourGapSum\ 件、$K=3$ 缺 \mtQfourMinPatchKThree\ 件，类别明细见 8.3 节的缺口表；中继无人机的补配不应按「够用即可」只补到推荐分区的水平，而至少要按下界 $K-1=\mtQfourRelayNeed$ 架预留。补配政策给两个参考阈值：$K=2$ 时把八类资源逐类各加 1 件，即有 \mtQfourThreshKTwoMOne\ 个分区可以落地；$K=3$ 时要逐类各加 2 件，才出现 \mtQfourThreshKThreeMTwo\ 个可落地的分区。

\textbf{第三，提高返航安全余量要付运力代价，且这份代价并不均匀。}$\rho$ 超过 $0.2992$ 先剥夺 A 型对 S008 的可达性，超过 $0.3513$ 则方案整体失效。因此加严安全余量不能单独做，必须与「增加架次或调整机型配比」同时进行，否则会在某个阈值处直接失去对个别服务区的可达性。

\textbf{第四，通信保障的投入方向是时序与站址，而不是能量。}本方案中继能量占用比最高仅 $0.740$，六站高度落在 $65\sim185$\,m、没有一站触及 $300$\,m 的上限，瓶颈在链长与站址几何。这一判断同样带条件：换站址、加大服务窗口或放宽并发都会改变该占用比，届时能量完全可能先触顶，故它只描述本方案。
